%%
%% frontmatter/preface.tex
%%
%% Preface: Why MatinBook, who is this book for, what is
%% its philosophy.
%%

\chapter*{پیشگفتار}
\addcontentsline{toc}{chapter}{پیشگفتار}
\label{ch:preface}

\section*{چرا \lr{MatinBook}؟}

حروف‌چینی کتاب‌های فنی فارسی، سال‌هاست با چالش‌های
متعددی روبه‌روست. ابزارهای موجود، یا برای زبان فارسی
بهینه نشده‌اند، یا از نظر بصری کیفیت لازم را ندارند،
یا برای پروژه‌های بزرگ مقیاس‌پذیر نیستند. این وضعیت،
باعث شده که بسیاری از کتاب‌های فنی فارسی، با وجود
محتوای ارزشمند، از نظر ظاهری فاصله‌ی زیادی با
استانداردهای نشر جهانی داشته باشند.

\lr{MatinBook} با هدف پر کردن این شکاف طراحی شده است.
این کلاس \lr{LaTeX}، تمام ابزارهای لازم برای تولید
یک کتاب فنی حرفه‌ای فارسی را در یک بسته‌ی منسجم فراهم
می‌کند: از حروف‌چینی راست‌به‌چپ و پشتیبانی کامل از
فونت‌های فارسی، تا محیط‌های علمی رنگ‌آمیزی‌شده،
نمودارهای پیچیده، و مدیریت حرفه‌ای منابع و نمایه.

\section*{این کتاب برای چه کسانی است؟}

این راهنما برای چهار گروه اصلی نوشته شده است:

\begin{itemize}
    \item \textbf{نویسندگان:} کسانی که می‌خواهند کتاب
    فنی یا علمی فارسی بنویسند و به دنبال ابزاری حرفه‌ای
    برای حروف‌چینی هستند.

    \item \textbf{پژوهشگران:} کسانی که می‌خواهند مقالات،
    پایان‌نامه‌ها یا گزارش‌های فنی خود را با کیفیت
    بالا منتشر کنند.

    \item \textbf{دانشجویان:} کسانی که می‌خواهند با
    ابزارهای حرفه‌ای حروف‌چینی آشنا شوند و مهارتی
    ارزشمند برای آینده‌ی حرفه‌ای خود کسب کنند.

    \item \textbf{مدرسان:} کسانی که می‌خواهند جزوه‌ها،
    کتاب‌های درسی یا مطالب آموزشی خود را با استاندارد
    بالایی آماده کنند.
\end{itemize}

\section*{فلسفه‌ی طراحی}

\lr{MatinBook} بر سه اصل بنیادین استوار است:

\subsection*{۱. وضوح بر تزئین}

هر عنصر بصری در \lr{MatinBook}، نقشی معنایی دارد.
رنگ‌ها، جعبه‌ها، و نمودارها، نه برای زیبایی صرف،
بلکه برای کمک به درک بهتر محتوا طراحی شده‌اند.
این اصل، از فلسفه‌ی نشر \lr{MIT Press} و
\lr{Springer} الهام گرفته است.

\subsection*{۲. سازگاری بر نوآوری}

هر عنصر در \lr{MatinBook}، یک شکل استاندارد دارد
و در تمام کتاب، یکسان ظاهر می‌شود. این سازگاری،
باعث می‌شود خواننده بدون حواس‌پرتی، روی محتوا
تمرکز کند.

\subsection*{۳. خویشتن‌داری بر افزونگی}

فضای خالی، یک عنصر طراحی است. \lr{MatinBook}
از پر کردن هر اینچ صفحه پرهیز می‌کند و به متن
اجازه می‌دهد نفس بکشد. این اصل، خوانایی را
افزایش می‌دهد و خستگی چشم را کاهش می‌دهد.

\section*{ساختار این راهنما}

این راهنما در شش بخش اصلی و پنج پیوست سازماندهی
شده است:

\begin{enumerate}
    \item \textbf{شروع به کار:} معرفی، نصب، اولین
    کتاب، ساختار کلی.

    \item \textbf{اجزای کتاب:} جلد، فصل، متن، محیط‌های
    علمی، ریاضیات، کد، نمودار، جدول، یادداشت حاشیه،
    منابع، نمایه، ارجاع.

    \item \textbf{شخصی‌سازی:} آپشن‌ها، رنگ‌ها، فونت‌ها،
    صفحه‌آرایی، تم‌ها، زبان.

    \item \textbf{مباحث پیشرفته:} ماژول‌ها، تم سفارشی،
    دستورات سفارشی، پروژه‌های بزرگ.

    \item \textbf{مثال‌های کامل:} سه فصل نمونه از
    کتاب‌های ریاضی، برنامه‌نویسی و فیزیک.

    \item \textbf{مرجع دستورات:} مرجع کامل کلاس،
    محیط‌ها و دستورات.
\end{enumerate}

هر فصل با یک مثال ساده شروع می‌شود، به تدریج پیچیده‌تر
می‌شود، و با تمرین‌های عملی به پایان می‌رسد. خواننده
می‌تواند فصل‌ها را به ترتیب بخواند، یا هر فصل را
به‌عنوان یک مرجع مستقل استفاده کند.

\section*{سپاسگزاری}

\lr{MatinBook} بدون پشتیبانی جامعه‌ی \lr{LaTeX} فارسی
و تلاش‌های بی‌دریغ توسعه‌دهندگان بسته‌های متن‌باز،
امکان‌پذیر نبود. از \lr{Vafa Khalighi} برای
\lr{XePersian}، از \lr{Hassan Mesgarha} برای
\lr{XePersian-HM}، از \lr{Thomas F. Sturm} برای
\lr{tcolorbox}، از \lr{Geoffrey Poore} برای
\lr{minted}، و از تمام کسانی که در این مسیر
کمک کرده‌اند، صمیمانه سپاسگزارم.

\vspace{1cm}

\begin{flushleft}
    \textbf{متین محمودی} \\
    مهر ۱۴۰۵
\end{flushleft}

\clearpage